\subsection{Политика слияния}

LSM-деревья различаются между собой по политике слияния дисковых компонент. \cite{luo2020lsm} На практике обычно используют следующие две политики, каждая заточена под свой преобладающий тип запросов:

\begin{paracol}{2}
  \begin{center}\bfseries Leveling Merge Policy\end{center}

  Данная политика хранит компоненты следующим образом:

  На каждом уровне будет хранится ровно одна компонента.

  Максимальный объем компоненты на уровне $L$ хранит в $T$ раз больше ключей, чем компонента на уровне $L - 1$. Тогда компонента на слое $L$ сольется с компонентой на $L - 1$  слое много раз (в точности $T$) перед тем, как переполнится и начнёт сливаться с компонентой на слое $L + 1$.

  Как итог, Leveling Merge Policy хранит "мало" компонент (по сравнению со второй политикой), благодаря чему, показывает лучший перфоманс на запросах чтения. 

  
\switchcolumn
  \begin{center}\bfseries Tiering Merge Policy\end{center}

  Пусть вместимость $L-1$ слоя $P$ - ключей, тогда
  эта политика, на каждом слое $L - 1$ будет хранится не более $T$ компонент, каждая размера $\frac{P}{T}$ с возможно пересекающимся диапазоном ключей. Когда на слое наберётся $T$ заполненных компонент, сольем их в одну компоненту и положим на следующий уровень. Действительно, такое поведение будет корректным, потому что компонента на $L$ слое вместит в себя $T$ раз больше ключей, т.е. $P$, и максимальная вместительность уровня $L$ в точности $P \cdot T$, что в $T$ раз больше, чем на предыдущем уровне. И заполненный $L-1$ уровень станет ровно одной компонентой на $L$ слое.

  Эта политика, напротив, оптимизирована для запросов записи, за счёт снижения числа тяжёлых слияний, когда структуре приходится сливать старшие уровни.


\end{paracol}




Вообще говоря, классической политикой для LSM-дерева является именно Leveling Merge Policy, но параллельно LSM-дереву,  Jagadish и др. в своей статье \cite{jagadish1997incremental} представили очень похожую на LSMT структуру, оптимизированую под запросы записи, в которой применили "stepped-merge policy". Которая в последствии превратилась в Tiering Merge Policy, используемую в современных СУБД.
